% AnchorMTM with updated correspondence and Storms teaser, 2026-11-02.
% Build: tectonic -X compile --keep-logs method_anchor_figures_updated_en.tex
% Unchanged companion: method_anchor_single_parameter_supplement.tex/.pdf.
% Changes: ANCHORMTM_SINGLE_PARAMETER_VALIDATION.md.
\documentclass[review]{vgtc}
\usepackage[T1]{fontenc}
\usepackage{mathptmx}
\usepackage{cite}
\usepackage{amsmath,amssymb,amsthm}
\usepackage{placeins}
\usepackage{float,stfloats,needspace}
\usepackage{booktabs,multirow,colortbl}
\makeatletter
\setlength{\@dblfptop}{0pt}
\setlength{\@dblfpsep}{12pt}
\setlength{\@dblfpbot}{0pt plus 1fil}
\makeatother
\usepackage[ruled,vlined]{algorithm2e}
\SetKwInput{KwData}{Data}
\SetKwInput{KwResult}{Result}
\SetKwProg{Fn}{Function}{:}{}
\SetKw{KwReturn}{return}
\SetKw{KwBreak}{break}
\SetKwFunction{Construct}{ConstructTemporalMap}
\SetKwFunction{Refine}{RefineOrder}
\SetKwFunction{Barriers}{OptimalBarriers}
\SetAlFnt{\small}
\SetAlCapFnt{\small\bfseries}
\SetAlCapNameFnt{\small\bfseries}
\SetAlCapSkip{3pt}
\DontPrintSemicolon
\setcounter{secnumdepth}{3}
\newtheorem{theorem}{Theorem}
\newtheorem{lemma}{Lemma}
\newtheorem{proposition}{Proposition}
\newcommand{\dtop}{d_{\mathrm{top}}}
\newcommand{\lca}{\operatorname{lca}}
\newcommand{\nroe}{\mathrm{NROE}}
\newcommand{\pos}[1]{\left[#1\right]^+}
\title{AnchorMTM: Spatially Anchored Merge Tree Maps\\for Time-Varying Scalar Fields}
\author{Anonymous Authors}
\onlineid{0}
\vgtccategory{Research}
\vgtcinsertpkg
\hypersetup{pdftitle={AnchorMTM: Spatially Anchored Merge Tree Maps for Time-Varying Scalar Fields},pdfauthor={Anonymous Authors}}
\abstract{Temporal merge tree maps summarize evolving scalar hierarchies, but their display axis does not identify a fixed spatial location. Giving that axis a spatial reference creates a conflict: an exact profile keeps each subtree contiguous, while projected features can interleave. We derive three layout-independent lower bounds that quantify the displacement forced by exact topology, finite feature widths and a bounded merge-level error. AnchorMTM approaches these bounds through topology-budgeted leaf ordering, capped anchor placement and scalar-profile reconstruction. Controlled four-leaf experiments recover the bounds by exhaustive ordering and independent linear programs. Comparisons on five planar and two volumetric sequences show how anchoring improves reference alignment and how a declared merge allowance reduces displacement. Recomputed feature-position and per-frame distribution plots locate these changes across complete sequences: a small budget resolves most Cylinder conflicts, whereas Tangaroa requires larger allowances.}
\keywords{Time-varying scalar fields, merge tree maps, spatial anchoring, topology preservation, constrained layout.}
\teaser{
\centering
\includegraphics[width=\linewidth]{figures/anchor_storms_overview.pdf}
\par
\refstepcounter{figure}\label{fig:storms_overview}
{\captionfonts\raggedright\noindent\figurename~\thefigure: Storms temporal overview ($\varepsilon=5\%$), using the case-study preprocessing. Spatial snapshots and profiles identify one matched extremum; blue outlines show leaf-arc supports. Rotated profiles and dashed intermediate samples connect to its trajectory in the full 744-frame map.\par}

}
\linespread{1.0}
\setlength{\textfloatsep}{7pt}
\setlength{\floatsep}{6pt}
\setlength{\dbltextfloatsep}{7pt}
\setcounter{dbltopnumber}{2}
\renewcommand{\dbltopfraction}{0.9}
\renewcommand{\dblfloatpagefraction}{0.7}
\renewcommand{\floatpagefraction}{0.7}
\setcounter{totalnumber}{6}
\setcounter{topnumber}{3}
\setcounter{bottomnumber}{3}
\renewcommand{\textfraction}{0.07}

\newcommand{\paperfigure}[3]{\centering\includegraphics[width=\textwidth]{#1}\caption{#3}\label{#2}}

\begin{document}
\firstsection{Introduction}
\maketitle
\raggedbottom
\setlength{\emergencystretch}{1.5em}
An analyst following storms, pollution peaks or vortices needs to see which features exist, where they are and how they move. Animation retains the spatial field at each instant; a static temporal overview places distant times side by side~\cite{Tversky:2002:Animation,Robertson:2008:Animation,Bach:2017:STC}. When time occupies one display axis, the coordinate assigned to the other determines which spatial relationships remain readable. Figure~\ref{fig:storms_overview} connects source-field supports to the trajectories of a merge tree map.

Temporal Merge Tree Maps (TMTM) compress each scalar field into a one-dimensional profile: extrema become minima and intervening barriers encode merge levels~\cite{Koepp:2023:TMTM}. Stacking the profiles preserves the evolving hierarchy. Spatiotemporal Merge Tree Maps (ST-MTM) further fit spatial distances and temporal position constraints~\cite{STMTM:2026}. Their coordinate expresses a tree traversal or relative geometry. Reading an anchor as longitude, or its displacement as eastward motion, requires an additional fixed spatial reference.

A spatial reference encounters a structural obstruction. An exact profile keeps every strict subtree contiguous, but projected sibling subtrees can interleave. If leaves $A,B$ belong to one subtree and $C,D$ to another, the reference order $A,C,D,B$ splits the first block. Some features must leave their targets, or some displayed merge levels must change. Finite feature widths add packing costs. The resulting displacement has two causes: a part imposed by the hierarchy and a part introduced by the layout.

We quantify this distinction before constructing a map. Three bounds measure the minimum displacement forced by an exact point order, by finite-width intervals and by the restrictions surviving a merge-error allowance. The point bound is attained by recursive midpoint ordering; the interval bounds provide certificates against which a chosen order can be assessed. They depend on the tree, reference and widths, so the same analysis applies across layout objectives.

AnchorMTM uses these quantities to give the temporal map a shared spatial axis. It fixes a direction from matched feature motion or an analyst's choice, selects a leaf order within a merge budget, places anchors under a displacement cap and reconstructs the scalar profiles. Experiments follow the same argument: test the bounds, measure reference positioning, then inspect what additional merge error buys in complete sequences.

Our contributions are:
\begin{itemize}\setlength{\itemsep}{1pt}\setlength{\parsep}{0pt}
\item Three layout-independent position bounds for exact topology, finite widths and bounded merge error.
\item A construction that couples fixed spatial targets with budgeted ordering and capped temporal placement.
\item Controlled-tree validation and planar/volumetric evidence connecting reference displacement, merge budgets and feature trajectories.
\end{itemize}

\section{Related Work}\label{sec:related}
\subsection{Spatial Coordinates in Temporal Overviews}
Static views and animation distribute temporal comparison across space and time~\cite{Tversky:2002:Animation,Robertson:2008:Animation}; generalized space--time cubes describe transformations between these organizations~\cite{Bach:2017:STC}. Hovm\"oller diagrams retain a physical coordinate while aggregating the remaining dimensions~\cite{Hovmoller:1949}. Geographic orderings and space-filling curves traverse the spatial domain~\cite{Guo:2006:SpatialOrder,Dafner:2000:CSFC}, with extensions for volume comparison and scalar-aware traversal~\cite{Weissenboeck:2019:DVL,Zhou:2021:DSF}. Their coordinate records traversal rank, while averaging can merge nearby extrema.

One-dimensional projections and MotionRugs summarize moving entities~\cite{Franke:2021:1DP,Buchmueller:2019:MR}. More generally, embeddings trade original-domain relationships against display organization~\cite{Nonato:2019:MP}. Figure~\ref{fig:related-work-mapping} contrasts sample traversal, selected features and hierarchy-preserving profiles. AnchorMTM assigns projected spatial targets to extrema while retaining the profile's scalar and hierarchical representation.

\begin{figure}[!tb]
\centering
\includegraphics[width=\columnwidth]{figures/anchor_related_work_mapping.pdf}
\caption{Mapping strategies on a common 35-sample field: domain traversal; a 22-sample feature summary; and a topology-based profile retaining the complete illustrated hierarchy.}
\label{fig:related-work-mapping}
\end{figure}
\subsection{Topological Profiles and Constrained Layouts}
Topology organizes scalar structures by connectivity~\cite{Heine:2016:Survey}. Landscapes and one-dimensional profiles express these hierarchies geometrically~\cite{Weber:2007:TL,Oesterling:2013:TLP}, including geometry-aware and temporally coherent extensions~\cite{Beketayev:2012:GPTL,Herick:2020:TCTL}. TMTM aligns profiles across time~\cite{Koepp:2023:TMTM}; ST-MTM adds spatial-distance fitting and temporal constraints~\cite{STMTM:2026}. AnchorMTM reuses this profile foundation and gives its axis a fixed origin, direction and scale.

Optimal leaf ordering minimizes adjacent dissimilarity subject to a hierarchy~\cite{BarJoseph:2001:OLO}; target-order reordering studies related compatible permutations~\cite{Bulteau:2022:RT}. Separation-constrained layouts manage overlap and displacement~\cite{Dwyer:2006:IPSEP,Li:2019:SOI}. Our bounds state the displacement that these ordering and packing constraints force relative to fixed targets. Merge-tree comparison depends on feature labels and correspondences~\cite{Yan:2021:Survey}; labelled cophenetic/interleaving representations support our fixed-identity merge-error measure~\cite{Munch:2019:LINF,Gasparovic:2024:IID}. Ordered merge trees connect scalar curves with interleaving distances~\cite{Beurskens:2025:OMT}, and minimax ultrametric fitting provides a related setting for correcting merge levels~\cite{Farach:1995:RM}.

\section{Spatially Anchored Merge Tree Maps}\label{sec:formulation}\label{sec:background}
\subsection{Trees, Profiles and Feature Identity}
For a sampled field $f_t$, the merge tree $T_t$ records sublevel components created at minima and joined at saddles; negating the field handles maxima~\cite{Tierny:2018:TTK}. Extraction and persistence simplification precede layout. A leaf $i$ has extremum location $p_i^t$ and support $S_i^t$, comprising the samples on its leaf arc. Augmented internal arcs retain samples for filling the profile. Write $h_t(v)$ for a node's scalar level and $\mathcal{D}_t(v)$ for its descendant leaves. The fixed-label merge matrix is
\begin{equation}\label{eq:merge}
V_{ij}^t=h_t(\lca(i,j)),\quad i\ne j;\qquad V_{ii}^t=f_t(p_i^t).
\end{equation}
Equal-level merges form one multifurcation. A one-dimensional profile $\psi_t$ places the extrema in leaf order $\pi_t$; the largest profile value between anchors $i,j$ gives their displayed merge level $V_{ij}^{\psi_t}$. Exact topology requires each strict subtree to form one contiguous block.

On a fixed grid, one-to-one Hungarian matching maximizes $|S_i^{t-1}\cap S_j^t|$ and discards zero-overlap pairs~\cite{STMTM:2026}. A matched leaf inherits its track identity; $\mu_t(i)$ denotes its predecessor. These links supply the temporal targets and motion samples below. The support is the leaf arc rather than the full component: internal arcs keep their own samples for reconstruction. Matching maximizes intersection counts, so a continuation is an assigned computational correspondence; unmatched leaves create or end tracks. A birth therefore contributes a current reference target but no predecessor target. The same matching is reused for every tested budget, keeping the temporal information fixed while the order constraints change.

\begin{figure}[!tb]
\centering
\includegraphics[width=\columnwidth]{figures/anchor_node_mapping.pdf}
\caption{An augmented merge tree and its profile retain the identities of extrema and regular arc samples. Their ordering makes every strict subtree contiguous.}
\label{fig:node-mapping}
\end{figure}
\subsection{One Reference for the Whole Sequence}\label{sec:reference-model}
\begin{figure}[!tb]
\centering
\includegraphics[width=\linewidth]{figures/anchor_leaf_correspondence.pdf}
\caption{Matching leaf-arc supports by one-to-one maximum overlap continues tracks across frames; unmatched leaves begin or end tracks. These correspondences supply the reference-motion samples.}
\label{fig:leaf_correspondence}
\end{figure}
The default direction maximizes squared matched-extremum motion:
\begin{equation}\label{eq:autodirection}
a^*\in\arg\max_{\|a\|=1}a^\top\Bigl(\sum_{t,i}\Delta p_i^t(\Delta p_i^t)^\top\Bigr)a,
\end{equation}
where $\Delta p_i^t=p_i^t-p_{\mu_t(i)}^{t-1}$ over valid links. Its sign makes total projected motion nonnegative; the longest domain axis handles zero motion energy. An analyst may supply a physical direction instead. For nearly zero signed motion, the largest-magnitude direction component is made positive; repeated leading eigenvalues can be resolved using a chosen physical axis. This separates reference estimation from the layout search: a budget changes the admissible map, while the coordinate being approximated stays fixed. Projected domain corners define one origin $o$ and extent $L$, giving
\begin{equation}\label{eq:reference}
q_i^t=\langle a^*,p_i^t\rangle-o\in[0,L].
\end{equation}
Direction, origin and extent remain fixed across frames and budgets. Relative geometry uses support-centroid distances $d_{ij}^t$; the targets $q_i^t$ use extremum locations. Thus support growth can change relative-distance terms without moving a feature's spatial target.

\subsection{Layout Model and Errors}\label{sec:common}
Each leaf occupies an interval of width $W_i=s\,m_i$ and center $z_i$, where $m_i$ is its support measure and $s$ is fixed. The layout set $\mathcal{L}(\pi;W)$ places these intervals in order $\pi$ inside $[0,L]$, with gap at least $g$ and anchor slack $|x_i-z_i|\le\rho W_i/2$. Two order costs distinguish point ordering from packing:
\begin{equation}\label{eq:nroe}
\begin{aligned}
\nroe(\pi)&=\frac{\max_{j\le k}[q_{\pi(j)}-q_{\pi(k)}]^+}{2L},\\
\tau(\pi)&=\min_{(x,z)\in\mathcal{L}(\pi;W)}\|x-q\|_\infty .
\end{aligned}
\end{equation}
NROE is the least normalized worst displacement of points in order $\pi$; $\tau$ includes the intervals and is an exactly solvable minimax linear program. Merge error measures the displayed hierarchy:
\begin{equation}\label{eq:dtop}
\dtop(\psi,T)=\max_{i,j}|V_{ij}^{\psi}-V_{ij}|\le\eta.
\end{equation}
The diagonal includes extremum values. Under complete fixed labelling this cophenetic error equals labelled interleaving distance and bounds the unlabelled interleaving and bottleneck distances~\cite{Munch:2019:LINF,Gasparovic:2024:IID,Morozov:2013:IDMT}.

These variables separate layout priorities in the common model
\begin{equation}\label{eq:family}
\begin{aligned}
\min_{\pi\in\Pi,\,(x,z)\in\mathcal{L}(\pi;W)}\;&\lambda_RE_R+\lambda_GE_G+\lambda_ME_M+\lambda_CE_C\\
\mathrm{s.t.}\;&\dtop\le\eta.
\end{aligned}
\end{equation}
where the terms fit spatial references, support distances, matched temporal targets and interval centers. TMTM packs tree-compatible blocks and aligns them by overlap; ST-MTM fits relative geometry with temporal constraints. Both preserve the tree and have no absolute-reference penalty. Anchoring adds $E_R$. A projection can instead follow every $q_i$, but the order $A,C,D,B$ breaks the $A,B$ subtree block. This conflict motivates bounds that hold across objective choices.

\section{The Positional Price of Keeping the Tree}\label{sec:bounds}
We first bound the point-order obstruction, then include widths, then identify which constraints survive a merge allowance. Suppress $t$, write $[v]^+=\max(v,0)$ and take empty maxima as zero. The supplementary material gives the complete packing derivations; here each bound identifies the particular constraint that the construction must address.

\subsection{Exact Topology}
Let $[m_X,M_X]$ be the reference hull of subtree $X$ in the reduced equal-level tree.
\begin{theorem}[Exact-topology position bound]\label{thm:order}
Every exact profile satisfies $\nroe(\pi)\ge\nu^*$, where
\begin{equation}\label{eq:nu}
\nu^*=\frac{1}{2L}\max_v\max_{X\ne Y\in\mathrm{ch}(v)}[\min(M_X-m_Y,M_Y-m_X)]^+.
\end{equation}
Recursive ordering by child-hull midpoint attains $\nu^*$ in $O(n\log n)$ time.
\end{theorem}
Sibling blocks must take one orientation. Its inversion cost is either $[M_X-m_Y]^+/2$ or $[M_Y-m_X]^+/2$, and midpoint order chooses the smaller. Every inverted leaf pair separates at its LCA. Therefore error above $\nu^*$ can be removed by reordering in the point model.

\subsection{Finite Widths}
Set $H_i=\rho W_i/2$, $\ell_i=q_i-W_i/2-H_i$, $h_i=q_i+W_i/2+H_i$ and $s_i=W_i+g$. Packing summaries
\begin{equation}\label{eq:packing}
\begin{split}
A(S)&=\max_{a\in S}\bigl(\ell_a+\sum_{i\in S:\ell_i\ge\ell_a}s_i\bigr),\\
B(S)&=\max_{b\in S}\bigl(\sum_{i\in S:h_i\le h_b}s_i-h_b\bigr)
\end{split}
\end{equation}
measure the interval span forced by upper and lower reference-endpoint sets. Define $\omega(X,Y)=\tfrac12\min\{A(X)+B(Y),A(Y)+B(X)\}-g/2$, and $E_\partial=\max\{0,B(\mathcal F)-g,A(\mathcal F)-g-L\}$ for all leaves $\mathcal F$.
\begin{theorem}[Width-aware position bound]\label{thm:width}
For every feasible exact-topology interval layout,
\begin{equation}\label{eq:widthbound}
\|x-q\|_\infty\ge\tau(\pi)\ge\mathcal B_w
=\max\{E_\partial,\max_{(X,Y)\in\mathcal P_0}\omega(X,Y)\},
\end{equation}
where $\mathcal P_0$ contains distinct sibling descendant-leaf sets. With zero widths and gaps, $\mathcal B_w=L\nu^*$.
\end{theorem}
The first and last anchors enclose their intervals and gaps up to the permitted anchor slack; sibling contiguity fixes a block orientation. The two canvas ends supply $E_\partial$. The bound has coordinate units, while $\nu^*$ is normalized by $L$.

\subsection{A Bounded Merge Error}
If leaf $k$ lies between descendants $i,j$ of node $u$, interval containment gives $V_{ik}^{\psi}\le V_{ij}^{\psi}$, hence $V_{ik}\le h(u)+2\eta$. Let $A_\eta(u)$ be the highest ancestor with $h(A_\eta(u))-h(u)\le2\eta$. Form $\mathcal P_\eta$ from pairs $(\mathcal D(u),\{k\})$ with $k\notin\mathcal D(A_\eta(u))$, and pairs $(\mathcal D(u),\mathcal D(w))$ with each set disjoint from the other's $A_\eta$ ancestor set, for internal $u,w$.
\begin{theorem}[Bounded-topology position bound]\label{thm:chain}
If $\dtop\le\eta$, every pair in $\mathcal P_\eta$ lies wholly in one block orientation. Thus
\begin{equation}\label{eq:chainbounds}
\begin{split}
\nroe(\pi)&\ge C_\eta=\max_{(X,Y)\in\mathcal P_\eta}\frac{[\min(M_X-m_Y,M_Y-m_X)]^+}{2L},\\
\tau(\pi)&\ge C_{w,\eta}=\max\{E_\partial,\max_{(X,Y)\in\mathcal P_\eta}\omega(X,Y)\}.
\end{split}
\end{equation}
\end{theorem}
The accumulated ancestor gap determines when an outside leaf can enter a block. Flattening parent gaps of at most $2\eta$ gives a second certificate; we use the larger of the two valid bounds. These certificates describe an unavoidable cost. The exact fixed-order value $\tau$ then measures how closely a candidate approaches it.



\begin{figure*}[!tb]
\paperfigure{figures/anchor_framework.pdf}{fig:framework}{AnchorMTM: extract and match features; fix spatial references; search admissible leaf orders; optimize anchors under a displacement cap; reconstruct scalar profiles and stack them along time.}
\end{figure*}


\section{AnchorMTM: Approaching the Bounds}\label{sec:method}\label{sec:construction}

AnchorMTM follows the bounds in two stages: first select a legal order with low $\tau$, then balance spatial and temporal quality within a cap on that order's displacement. The control is a normalized merge allowance $\varepsilon$, with $\eta=(\varepsilon_0+\varepsilon)\Delta_F$ for sequence scalar range $\Delta_F$ and numerical tolerance $\varepsilon_0=10^{-7}$. Exact Mode sets $\varepsilon=0$; the exact experiments use zero labelled error. Widths, gap, slack and weights stay fixed during a budget sweep.

\subsection{Order Selection and Optimal Barriers}
A finite candidate pool $\mathcal C$ starts from midpoint orders of the strict tree, threshold flattenings and selective conflict flattenings. Width-aware block moves refine promising structures. Each visited order is evaluated against the original merge matrix, and selection is lexicographic:
\begin{equation}\label{eq:orderobjective}
\pi^\varepsilon\in\arg\min_{\pi\in\mathcal C}\bigl(\tau(\pi),\nroe(\pi)\bigr)
\quad\mathrm{s.t.}\quad\delta^*(\pi)\le\eta.
\end{equation}
The exact tree supplies a feasible strict-topology seed. Selective flattening releases additional permutations while retaining deeper blocks, as illustrated in Fig.~\ref{fig:selective_relaxation}.

\paragraph{Construction sequence.} For each frame, generate and refine candidate orders; compute $\tau$, NROE and $\delta^*$; admit the budget-feasible orders; choose the lexicographic minimum; solve its capped QP; fill and correct the profile; then check the rendered merge matrix and append the temporal column. The supplementary algorithm lists these steps with the tracking and rounding procedures.


\begin{figure}[!tb]
\centering
\includegraphics[width=\columnwidth]{figures/anchor_relaxation_open.pdf}
\caption{Selective flattening admits order $A,C,D,B$ while retaining the $C,D$ grouping. The displayed barrier levels are corrected against the original merge matrix.}
\label{fig:selective_relaxation}
\end{figure}
For gap $k$ of an order, set $C_k=\min_{i\le k<j}V_{\pi(i)\pi(j)}$ and let $F_k$ be the larger adjacent extremum value.
\begin{proposition}[Optimal barriers]\label{prop:barriers}
With fixed extrema, the least continuous merge error is
\begin{equation}\label{eq:delta}
\delta^*(\pi)=\max\Bigl\{0,\max_k(F_k-C_k),
\tfrac12\max_{i<j}\bigl(V_{\pi(i)\pi(j)}-\max_{i\le k<j}C_k\bigr)\Bigr\},
\end{equation}
attained by $b_k=C_k+\delta^*(\pi)$.
\end{proposition}
Indeed, error $\delta$ requires $F_k\le b_k\le C_k+\delta$; choosing the largest permitted barriers meets every pair's lower requirement exactly when $\delta\ge\delta^*$. This gives a closed-form admission test for candidate orders.

\subsection{Capped Spatiotemporal Placement}
For the chosen order, the default weights are $(\lambda_R,\lambda_G,\lambda_M,\lambda_C)=(4,1,0.5,0.1)$. Each term in Eq.~\ref{eq:family} is a mean squared residual: anchor-to-reference distance, ordered anchor separation minus support-centroid distance, matched temporal-target error, or anchor-to-center offset. The temporal target follows projected motion:
\begin{equation}\label{eq:temporal_target}
\hat x_i^t=x_{\mu_t(i)}^{t-1}+q_i^t-q_{\mu_t(i)}^{t-1}.
\end{equation}
Unmatched leaves have no temporal penalty. Placement solves the convex quadratic objective under the interval constraints and
\begin{equation}\label{eq:qp}
\|x-q\|_\infty\le\tau(\pi^\varepsilon)+\beta_0,
\end{equation}
with fixed slack $\beta_0$. Because $\tau$ is attained, this cap is feasible and transfers the selected order's displacement control to the final anchors. The cap prevents a relative-distance or temporal term from moving an anchor arbitrarily far from its spatial target. Conversely, the fixed slack leaves room to improve these secondary relationships after the discrete order has been chosen. Within a frame the feasible set is linear and the objective quadratic; frame order enters only through matched predecessor targets.

\subsection{Profile Reconstruction}
Leaf-arc samples fill the intervals and LCA paths fill their gaps~\cite{Koepp:2023:TMTM,STMTM:2026}. Barrier correction clips each inter-anchor interior to $b_k$ and raises its peak when needed. Profiles share direction, origin, extent and pixel scale; their stack is $I^\varepsilon(t,u)=\psi_t^\varepsilon(u)$. Integer rounding can affect merge levels, so the rendered fixed-label matrix, including its diagonal, is checked after reconstruction. The supplementary material supplies the full algorithm and rounding procedure.





\begin{table*}[!tb]
\caption{Five-dataset comparison. Bold/underline mark best/second distinct values among comparable local measurements only; ties share emphasis. Daggered SNS/TW values are published context and excluded from ranking. Curves use fixed default orientations.}
\label{tab:quantitative-final}
\footnotesize
\setlength{\tabcolsep}{1.5pt}
\renewcommand{\arraystretch}{0.85}
\begin{tabular*}{\textwidth}{@{\extracolsep{\fill}}l*{12}{c}lccc@{}}
\toprule
\multirow{2}{*}{Dataset} & \multicolumn{3}{c}{Hilbert} & \multicolumn{3}{c}{Morton} & \multicolumn{3}{c}{TMTM} & \multicolumn{3}{c}{ST-MTM} & \multicolumn{4}{c}{AnchorMTM (ours)} \\
\cmidrule(lr){2-4}\cmidrule(lr){5-7}\cmidrule(lr){8-10}\cmidrule(lr){11-13}\cmidrule(lr){14-17}
 & SNS$\downarrow$ & TW$\uparrow$ & NPE$\downarrow$ & SNS$\downarrow$ & TW$\uparrow$ & NPE$\downarrow$ & SNS$\downarrow$ & TW$\uparrow$ & NPE$\downarrow$ & SNS$\downarrow$ & TW$\uparrow$ & NPE$\downarrow$ & Mode & SNS$\downarrow$ & TW$\uparrow$ & NPE$\downarrow$ \\
\midrule
\multirow{2}{*}{Gaussian} & \multirow{2}{*}{$0.1117$} & \multirow{2}{*}{$0.6719$} & \multirow{2}{*}{$0.7557$} & \multirow{2}{*}{$\underline{0.0393}$} & \multirow{2}{*}{$\underline{0.8333}$} & \multirow{2}{*}{$0.7359$} & \multirow{2}{*}{$0.25^{\dagger}$} & \multirow{2}{*}{$0.77^{\dagger}$} & \multirow{2}{*}{$\underline{0.1118}$} & \multirow{2}{*}{$0.06^{\dagger}$} & \multirow{2}{*}{$1.00^{\dagger}$} & \multirow{2}{*}{$0.1154$} & Exact & $\mathbf{0.0353}$ & $\mathbf{0.9323}$ & $\mathbf{0.0089}$ \\
 &  &  &  &  &  &  &  &  &  &  &  &  & Balanced & $\mathbf{0.0353}$ & $\mathbf{0.9323}$ & $\mathbf{0.0089}$ \\
\multirow{2}{*}{Ring} & \multirow{2}{*}{$\mathbf{0.1252}$} & \multirow{2}{*}{$\mathbf{0.9048}$} & \multirow{2}{*}{$0.5941$} & \multirow{2}{*}{$0.1489$} & \multirow{2}{*}{$0.7619$} & \multirow{2}{*}{$0.6033$} & \multirow{2}{*}{$0.28^{\dagger}$} & \multirow{2}{*}{$0.68^{\dagger}$} & \multirow{2}{*}{$0.1756$} & \multirow{2}{*}{$0.21^{\dagger}$} & \multirow{2}{*}{$0.88^{\dagger}$} & \multirow{2}{*}{$0.2463$} & Exact & $0.1485$ & $\underline{0.8571}$ & $\underline{0.0471}$ \\
 &  &  &  &  &  &  &  &  &  &  &  &  & Balanced & $\underline{0.1334}$ & $\underline{0.8571}$ & $\mathbf{0.0395}$ \\
\multirow{2}{*}{Wildfires} & \multirow{2}{*}{$0.1883$} & \multirow{2}{*}{$0.8303$} & \multirow{2}{*}{$0.3641$} & \multirow{2}{*}{$0.1304$} & \multirow{2}{*}{$0.8758$} & \multirow{2}{*}{$0.1560$} & \multirow{2}{*}{$0.44^{\dagger}$} & \multirow{2}{*}{$0.88^{\dagger}$} & \multirow{2}{*}{$0.2280$} & \multirow{2}{*}{$0.24^{\dagger}$} & \multirow{2}{*}{$0.91^{\dagger}$} & \multirow{2}{*}{$0.3180$} & Exact & $\mathbf{0.0426}$ & $\mathbf{0.9136}$ & $\underline{0.0131}$ \\
 &  &  &  &  &  &  &  &  &  &  &  &  & Balanced & $\underline{0.0447}$ & $\underline{0.8989}$ & $\mathbf{0.0116}$ \\
\multirow{2}{*}{Storms} & \multirow{2}{*}{$0.1290$} & \multirow{2}{*}{$0.7403$} & \multirow{2}{*}{$0.3705$} & \multirow{2}{*}{$0.1180$} & \multirow{2}{*}{$0.7251$} & \multirow{2}{*}{$0.3468$} & \multirow{2}{*}{$0.34^{\dagger}$} & \multirow{2}{*}{$0.71^{\dagger}$} & \multirow{2}{*}{$0.5753$} & \multirow{2}{*}{$0.21^{\dagger}$} & \multirow{2}{*}{$0.94^{\dagger}$} & \multirow{2}{*}{$0.4579$} & Exact & $\underline{0.0790}$ & $\underline{0.8463}$ & $\underline{0.0399}$ \\
 &  &  &  &  &  &  &  &  &  &  &  &  & Balanced & $\mathbf{0.0743}$ & $\mathbf{0.8766}$ & $\mathbf{0.0140}$ \\
\multirow{2}{*}{PM2.5} & \multirow{2}{*}{$0.2386$} & \multirow{2}{*}{$\mathbf{0.8887}$} & \multirow{2}{*}{$0.8090$} & \multirow{2}{*}{$\underline{0.1760}$} & \multirow{2}{*}{$\underline{0.8277}$} & \multirow{2}{*}{$0.8823$} & \multirow{2}{*}{$0.3018$} & \multirow{2}{*}{$0.7781$} & \multirow{2}{*}{$0.5241$} & \multirow{2}{*}{$0.1897$} & \multirow{2}{*}{$0.8224$} & \multirow{2}{*}{$0.5577$} & Exact & $0.1964$ & $0.7884$ & $\underline{0.1390}$ \\
 &  &  &  &  &  &  &  &  &  &  &  &  & Balanced & $\mathbf{0.1636}$ & $0.7657$ & $\mathbf{0.0468}$ \\
\bottomrule
\end{tabular*}
\vspace{2pt}
\begin{minipage}{\textwidth}\footnotesize For the original four inputs, TMTM/ST-MTM Temporal SNS/TW ($\dagger$) are published values; other entries are local measurements. Author inputs and coverage are unconfirmed. Local TW coverage (Gaussian/Ring/Wildfires/Storms) is 64/64, 1/40, 4/32, 11/744; ineligible frames are omitted. Published TW has two-decimal precision. Mode applies only to AnchorMTM. PM2.5 uses 744 shared-input local frames (TW $k=1$, full coverage); its Balanced mode uses the fixed exploratory 5\% case budget, without validation-set budget selection. ST-MTM uses the Storms preset with start frame 0; PM2.5 baseline metrics are all recomputed.\end{minipage}

\end{table*}


\begin{figure*}[!tb]
\paperfigure{figures/volume_gallery.pdf}{fig:volume-gallery}{Recomputed feature positions on Tangaroa and Cylinder using the default-weight volume configuration. Grey reference points and signed-displacement colors connect feature positions to $q$. Each panel reports the feature-pooled mean absolute displacement and NPE (mean per-frame maximum); positive budgets also report maximum attained merge error. Baseline reflection minimizes pooled reference displacement. These runs use the fixed configurations described in the protocol, distinct from the selected geometry-oriented table settings.}
\end{figure*}
\section{Experiments}\label{sec:experiments}
The experiments ask three connected questions: do the position bounds match controlled optima; how much does anchoring improve reference alignment; and where does a merge budget reduce displacement? The same fixed-label error and reference coordinate connect theory, tabulated results and per-frame plots.


\subsection{Data, Metrics and Configurations}
Five planar inputs are Gaussian (64 frames), Ring (40), Wildfires (32), Storms (744) and hourly PM2.5 (744). The first four use per-frame persistence fractions $(0.001,0.001,0.15,0.105)$ following the ST-MTM benchmark~\cite{STMTM:2026}. They are local reconstructions: three Gaussian peaks, the archived Ring generator, aggregated wildfire radiative power and ERA5 pressure anomalies. PM2.5 interpolates the January 2017 CAQRA hourly fields to a $160\times144$ grid, with 5\% per-frame persistence. Raster lengths are 4096, except Ring's original 196 positions.

The volumes are Cylinder (102 cropped $135\times64\times48$ $Q$-criterion fields, 2\% per-frame persistence)~\cite{Camarri:2005:Cylinder,Weinkauf:SquareCylinder} and Tangaroa (201 $300\times180\times120$ speed fields, range $[0,2]$, absolute persistence 0.2)~\cite{Popinet:2004:Tangaroa}. Within each comparison, methods share fields, retained labels and support geometry.

SNS measures scale-normalized support-distance stress~\cite{Smelser:2025:Scale}; TW measures neighborhood preservation~\cite{Venna:2001:NP}. Absolute positioning uses
\begin{equation}\label{eq:experimental-npe}
\mathrm{NPE}=\frac1T\sum_t\frac{\max_i|x_i^t-q_i^t|}{L}.
\end{equation}
For the original four planar inputs, published TMTM/ST-MTM SNS/TW provide context (daggers in Table~\ref{tab:quantitative-final}); NPE is measured locally. Local TW coverage is 64/64, 1/40, 4/32 and 11/744 for $k=(1,3,5,3)$. PM2.5 and volume scores are all local, with $k=1$ and full coverage. Hilbert/Morton use fixed orientations and full-grid ranks. Thus NPE compares a common reference, while SNS/TW evaluate complementary spatial relationships.

Exact uses zero allowance. On the four benchmark inputs, Balanced selects a positive budget from 1--10\% on the first chronological quarter by minimizing the equal-weight normalized NPE/SNS score, with smaller budgets breaking near ties; the resulting budgets are 1, 1, 5 and 10\%. PM2.5 retains its fixed 5\% mode. Sequence-level references and full-sequence reporting include these selection frames. Volume tables use fixed $+x$ references and selected geometry-oriented configurations; their complete parameters are in the supplement. The volume position/distribution figures use a second, common default-weight configuration, with motion-derived direction for Tangaroa and $x$ for Cylinder.

\begin{figure*}[!tb]
\paperfigure{figures/budget_distribution_complete.pdf}{fig:budget-distribution}{Recomputed per-frame budget response. Upper row: worst displacement, with quartiles, 5th--95th percentile whiskers, outliers and mean diamonds (square-root scale). Lower row: measured merge error and budget guides. Ring, Storms and PM2.5 use budgets 0, 1, 2, 5, 10 and 20\%; the two default-weight volumes use 0, 2, 5, 10 and 20\%. Orange marks the selected planar modes. Every plotted frame passes its continuous and raster budget checks.}
\end{figure*}
\subsection{Q1: Do the Bounds Recover the Forced Displacement?}\label{sec:bounds-eval}

\begin{figure}[!tb]
\centering
\includegraphics[width=\columnwidth]{figures/anchor_controlled_bounds.pdf}
\caption{Independent four-leaf validation. (a) The width-aware bound and order search coincide with exhaustive minimax LPs over 24 orders. (b) For four parent gaps, the displacement disappears when twice the merge allowance reaches that gap. Circles mark the recomputed bounds.}
\label{fig:controlled}
\end{figure}
Two sibling pairs merge below a root at 1, with leaf values 0 and targets $(0.2,0.5,0.5-d,0.8-d)$. Across 43 overlaps $d\in[-0.12,0.30]$, independent minimax LPs evaluate every leaf permutation. With width 0.06, gap 0.01 and $\rho=0.5$, the exhaustive optimum, width certificate and selected order coincide at $\tau=\max(0,(d+0.04)/2)$ (Fig.~\ref{fig:controlled}a). At $d=0.24$, the point cost is 12\% and the interval cost 14\%.

For parent gaps $h\in\{0.2,0.4,0.6,0.8\}$ and 91 budgets each, all 364 recomputed tests recover the release threshold $2\varepsilon=h$ (Fig.~\ref{fig:controlled}b). These tests connect the three bounds: interleaving causes the point obstruction, widths add packing, and scalar merge gaps determine the price of releasing a block.



\paragraph{Current data sequences.}
The same calculation applies to every frame of our five planar inputs. Table~\ref{tab:real-bounds} recomputes their bounds from the simplified source trees, extremum targets and fixed widths. The selected Exact order attains the width-aware bound on 1,335 of 1,624 frames. On those frames the certificate and an admissible solution coincide, establishing the optimum over all exact-topology orders. Gaussian and Ring attain it on every frame; Wildfires misses it on one and Storms on 15. PM2.5 attains it on 471 of 744 frames and leaves a larger residual gap: mean normalized $\tau$ is 13.11\%, versus a mean bound of 12.88\%. This is the interval-order cost before the capped placement objective chooses the final anchors.








\begin{figure*}[!tb]
\paperfigure{figures/anchor_storms_budget_maps.pdf}{fig:storms-budget-profiles}{Storms event window, frames 355--460, at five AnchorMTM budgets and with local TMTM/ST-MTM maps on the same case-study fields and trees. Common crops and the $[-35,35]$ hPa scale connect the map trajectories to the reference feature.}
\end{figure*}
\subsection{Q2: How Closely Do Maps Follow the Reference?}\label{sec:reference-results}
The baseline columns ask how much order displacement is forced by the point hierarchy. For each frame we sort its recorded anchors and evaluate NROE against the same $q$. The displayed share is $1-\sum_t\nu_t^*/\sum_t\nroe(\pi_t)$, so midpoint ordering removes that share without changing the tree in the point model. On Storms the excess is 91.0\% for TMTM and 83.0\% for ST-MTM; on PM2.5 it is 56.4\% and 59.7\%. The width certificate then identifies the packing cost that remains after order improvement. Together these measurements turn the distinction between forced and chosen displacement into a per-sequence diagnostic.
\begin{table}[!tb]
\centering\caption{Current planar sequences: Exact orders attaining the width-aware lower bound, and the share of local baseline point-order error above the exact point bound. Orientations match the primary measurements.}\label{tab:real-bounds}
\footnotesize\setlength{\tabcolsep}{3pt}
\begin{tabular}{@{}lrrcc@{}}\toprule
Dataset & Frames & Width hits & TMTM excess & ST-MTM excess\\\midrule
Gaussian & 64 & 64 & 100.0\% & 100.0\% \\
Ring & 40 & 40 & 64.5\% & 25.7\% \\
Wildfires & 32 & 31 & 92.5\% & 98.1\% \\
Storms & 744 & 729 & 91.0\% & 83.0\% \\
PM2.5 & 744 & 471 & 56.4\% & 59.7\% \\
\bottomrule\end{tabular}\end{table}

Balanced has lower NPE than the fixed curves and both local merge-tree adapters on all five planar inputs (Table~\ref{tab:quantitative-final}). Against curves, it also improves SNS and TW on Gaussian, Wildfires and Storms. Ring favors Hilbert's SNS/TW (0.1252/0.9048 versus 0.1334/0.8571), and PM2.5 favors Hilbert's TW (0.8887 versus 0.7657). The anchored axis improves location reading while retaining a different balance of local geometry.

\input{tables/table_volume.tex}
In the selected fixed-$x$ configurations, Balanced lowers SNS and NPE and raises TW relative to TMTM on both volumes (Table~\ref{tab:volume-comparison}). Cylinder NPE falls from 47.69\% to 4.58\%; Tangaroa falls from 38.66\% to 6.53\%. The selected Tangaroa configuration has zero reference weight and a 33.33\% allowance, so the mechanism is primarily geometry fitting within the prescribed reference cap.


Figure~\ref{fig:volume-gallery} reveals the spatial patterns behind reference error. Its default Exact configuration has NPE 11.18\% on Tangaroa and 9.18\% on Cylinder, versus TMTM's 30.63\% and 47.69\%. A 5\% budget brings Cylinder to 1.19\%, while Tangaroa reaches 8.91\% and needs 20\% to reach 1.88\%. The dark-grey reference columns now expose the trajectories against which these displacements are measured.








\begin{figure*}[!tb]
\paperfigure{figures/anchor_cylinder_centered_blue.pdf}{fig:cylinder-centered}{Cylinder: a central-tree guidance variant and five source $Q$ isosurfaces, using a common camera and blue scalar scale. Stages I--V identify early, transition, shedding, transition and late configurations. Its ordinate expresses the guided layout position.}
\end{figure*}
\subsection{Q3: What Does the Merge Budget Buy?}\label{sec:relaxation-results}
The primary Exact planar runs have zero fixed-leaf merge error. Balanced attains maximum normalized errors 0, 0.982\%, 4.421\% and 9.985\% on Gaussian, Ring, Wildfires and Storms, within their declared budgets. Relative to Exact, NPE decreases by 16.2\%, 11.6\% and 64.9\% on the latter three inputs; Gaussian remains unchanged. Wildfires SNS rises from 0.0426 to 0.0447, showing that improved anchoring can change relative geometry in the other direction.


Figure~\ref{fig:budget-distribution} localizes this improvement. On Storms, a 10\% allowance changes the median displacement from 1.10\% to 0.84\%, but lowers the 90th percentile from 10.7\% to 2.5\%; 482 of 744 frames retain zero merge error. The budget therefore acts chiefly on a difficult tail. In the default volumes, 200 of 201 Tangaroa frames and 98 of 102 Cylinder frames use a positive allowance, shifting most of the distribution.

The completed PM2.5 sweep uses the same input, tracks, reference, widths and weights at all six budgets. Its NPE decreases through 13.90\%, 10.30\%, 8.33\%, 4.68\%, 2.15\% and 1.15\%. All 2,976 newly solved frames at 1, 2, 10 and 20\% pass continuous and raster checks. TW changes from 0.7884 in Exact to 0.7580 at 20\%, linking the anchoring gain to a neighborhood cost.

\subsection{Reading the Resulting Maps}

\begin{figure}[!tb]
\centering
\includegraphics[width=\columnwidth]{figures/anchor_ring_result.pdf}
\caption{Ring: joint spatial guidance and merge allowances 0, 2, 5, 10 and 20\%, with same-input TMTM/ST-MTM maps, supplied curve panels and five source snapshots. All computed maps retain 40 frames and 196 positions.}
\label{fig:ring-result}
\end{figure}
\paragraph{Ring.} Figure~\ref{fig:ring-result} varies guidance jointly with budget: $q_{layout}=(1-\alpha)q_{mid}+\alpha q_{physical}$, with $\alpha=\varepsilon/20\%$ and a fixed $-128.7^\circ$ physical direction. Increasing guidance moves the ring toward its projected positions; the snapshots connect this change to expansion and strengthening.








\paragraph{Storms.} The case-study preprocessing of Fig.~\ref{fig:storms_overview} supplies the common fields for Fig.~\ref{fig:storms-budget-profiles}. The 0, 2 and 5\% maps overlap locally; 10 and 20\% move the selected feature toward its reference. TMTM aligns subtrees by overlap, while ST-MTM fits relative geometry with temporal constraints. The common event window lets the reader follow the resulting trajectory differences.

\begin{figure}[!tb]
\centering
\includegraphics[width=\columnwidth]{figures/PM25_single_column.pdf}
\caption{Hourly PM2.5, January 22 00:00--27 12:00, 2017 (Beijing time): the 5\% map and five source snapshots. All 133 hourly samples are retained; concentrations above 400 $\mu$g/m$^3$ saturate.}
\label{fig:pm25-hourly}
\end{figure}
\paragraph{PM2.5.} Figure~\ref{fig:pm25-hourly} places a 133-hour event within the complete January sequence. Join trees of negated concentration represent high-concentration features, and the display restores positive concentration. The 5\% mode attains 4.998\% maximum merge error across all 744 hours. Snapshots provide the geographic context for the matched peak visible in the temporal map.


\paragraph{Reproduction and cost.} The supplied code rebuilds the three added result figures from labelled layouts and raster columns: 43 controlled overlaps, 364 controlled budgets, 9,168 planar frame--budget observations and 1,515 volumetric columns. For each volume column it reconstructs adjacent barriers from the original scalar raster, accumulates pairwise maxima and compares the complete fixed-label matrix with the source tree. It also verifies support IDs, interval order, canvas capacity, gaps and anchor slack. Every chart mark is then generated from the resulting CSV records. The five-sequence audit separately recomputes the point/width certificates and the baseline inversion errors in Table~\ref{tab:real-bounds}. These two checks link the rendered spatial-error distributions to the theoretical quantities used during ordering. Planar optimization takes 0.15--3.32\,s on the four benchmark inputs and 199.97\,s for the uncached PM2.5 Balanced run on Apple M4; full pipeline times and configurations appear in the supplement.








\paragraph{Cylinder.}\label{sec:centered-volume} Central-tree guidance adds $16\,\mathrm{mean}_i(x_i-g_i)^2$ to the placement objective and reveals unfolding in Fig.~\ref{fig:cylinder-centered}. A 1\% allowance and 56.21\% position cap yield zero measured retained-leaf merge error. The snapshots at $t=32.1$, 56.1, 80.2, 120.3 and 162.0 link the five visual stages to the three-dimensional structures. This guidance changes the coordinate's role: frame-balanced absolute $X$ error rises from 21.1\% to 28.1\% as its weight grows from 1 to 16, trading location accuracy for hierarchical unfolding.

\paragraph{Configuration and runtime.} Tables~\ref{tab:anchor-parameters-final} and~\ref{tab:volume-parameters} give the fixed settings used by the quantitative comparisons. Only the merge allowance changes between each Exact--Balanced pair. The table runtimes measure one uncached run on Apple M4; optimization includes order search and anchor placement, while the total includes tree extraction, matching, filling and checks.
\input{tables/table2_parameters_runtime.tex}
\input{tables/table_volume_parameters.tex}
\Needspace{7\baselineskip}
\section{Discussion and Conclusion}\label{sec:discussion}
A spatial axis has a quantifiable hierarchical cost. The point bound identifies order error removable without changing the tree; finite-width and budgeted bounds identify remaining constraints. AnchorMTM follows this distinction by choosing an order before optimizing its placement. The experiments connect the controlled optima to real feature trajectories: budgets act on a small difficult tail in Storms and on most frames in the two default volumes.

Absolute position, relative geometry and hierarchical unfolding serve different readings of the same map. The curve comparisons and Cylinder guidance example make these choices visible. Tracking quality controls the estimated direction, and a single projection merges perpendicular spatial differences; source snapshots and analyst-selected directions address these ambiguities. Finite-width ordering remains heuristic, and a gap above its lower bound can originate in either the search or the certificate. These issues motivate tighter ordering methods and task-based studies of how readers use the anchored coordinate.

The result is a temporal merge tree map whose coordinate can be read against a fixed spatial reference, with an explicit merge allowance and measured displacement. Three bounds explain the restrictions; the reconstruction checks connect those restrictions to the displayed map.



\FloatBarrier
\begingroup
\interlinepenalty=10000
\bibliographystyle{abbrv-doi-narrow}
\bibliography{refs}
\endgroup
\end{document}
